\thispagestyle{plain}
\chapter*{
  \vspace{-1.7cm}
  \centering
  \rule{\textwidth}{0.5pt} \\[0.4cm]
  Conclusion et perspectives
}
\addcontentsline{toc}{chapter}{\numberline{}Conclusion et perspectives}

Au terme de ce projet de fin d'études, nous avons conçu et développé \textbf{Sellio}, une plateforme transactionnelle intelligente qui permet à des commerçants de huit secteurs d'activité de créer, personnaliser et exploiter leur propre boutique en ligne. Partie d'une boutique unique, NaturEssence, la solution est devenue une plateforme multi-boutiques fondée sur une architecture microservices : cinq services Spring Boot derrière une passerelle API, chacun propriétaire de sa base PostgreSQL, et deux applications React. \newline

La plateforme couvre l'ensemble du cycle d'une boutique : intégration et vérification des commerçants (KYC par carte bancaire, pièce d'identité et selfie), catalogue dont la fiche s'adapte au secteur et au type de chaque produit, huit modèles de vitrine dont l'apparence et la page d'accueil se personnalisent sans code avec un aperçu en direct, isolation des données de chaque boutique et de ses clients, tunnel de commande avec paiement Stripe, gestion des commandes, des retours, de la TVA et de la livraison, promotions, e-mailing et programme de fidélité. \newline

Elle se distingue par deux volets intelligents. Le premier est une \textbf{cabine d'essayage virtuel en temps réel}, fondée sur le modèle génératif \textit{Lucy VTON} de Decart et intégrée de façon sécurisée par des jetons éphémères. Le second est un ensemble de \textbf{modèles d'apprentissage automatique} alimentés par le suivi comportemental des visiteurs : un moteur de recommandation hybride, une segmentation des visiteurs par K-Means et une prédiction de l'attrition adaptée au secteur de chaque boutique. Ces modèles ont été entraînés et évalués sur trois jeux de données publics réels, nettoyés par des scripts reproductibles, avec un découpage temporel et une comparaison à des méthodes de référence. Les résultats expérimentaux montrent une amélioration relative de 10 à 16\,\% du HR@10 pour la recommandation, une AUC d'environ 0,80 pour l'attrition et des segments stables et interprétables. \newline

La dernière release a renforcé la qualité et la fiabilité de Sellio grâce à 230 tests automatisés couvrant notamment la sécurité, l'isolation des boutiques, les droits d'accès, les fonctionnalités métier et les interfaces. La plateforme est conteneurisée avec Docker et peut être lancée en une seule commande. Un pipeline GitHub Actions automatise la compilation, les tests et la construction des images Docker, sans publication sur un registre ni déploiement automatique sur un serveur. \newline

Sur le plan professionnel, ce projet nous a permis de consolider nos compétences en architecture distribuée, en sécurité des applications web, en intégration de services tiers (paiement, IA générative, e-mail), en ingénierie des données et en apprentissage automatique, tout en appliquant une démarche Agile fondée sur Scrum. Il nous a surtout appris l'importance d'une évaluation rigoureuse et mesurable pour juger la pertinence des solutions développées. \newline

\textbf{Perspectives.} Plusieurs évolutions peuvent être envisagées. Sur le plan des paiements, l'intégration d'un prestataire de paiement local prenant en charge le dinar tunisien, la confirmation des commandes par webhook et la revérification côté serveur des prix et du stock constituent des étapes indispensables avant une mise en production. La vérification des commerçants pourrait également être automatisée par une \textbf{reconnaissance faciale} comparant le selfie à la photo de la pièce d'identité, avec une intervention administrative pour les cas nécessitant une vérification supplémentaire. \newline

Sur le plan de l'architecture, une communication asynchrone par messages entre les microservices, l'hébergement de chaque base de données sur un serveur distinct, l'\textbf{orchestration des conteneurs} avec Kubernetes et la mise en place d'un \textbf{déploiement continu dans le Cloud}, incluant la publication des images Docker sur un registre, permettraient d'améliorer la disponibilité et la capacité de montée en charge. \newline

Sur le plan de l'intelligence, les modèles d'attrition pourraient être évalués sur d'autres jeux de données adaptés aux différents secteurs, comme Olist et Instacart, puis réentraînés sur les données réelles de Sellio lorsque les boutiques disposeront d'un historique suffisant. Les recommandations pourraient également être évaluées en production à travers des \textbf{tests A/B} afin de mesurer leur influence réelle sur les ventes. Enfin, le développement d'une \textbf{application mobile} destinée aux clients et aux commerçants permettrait d'élargir l'accessibilité de la plateforme. \newline

Ainsi, Sellio pourrait évoluer d'une plateforme de création de boutiques vers une véritable solution de commerce intelligent, locale, sécurisée et pilotée par les données de chaque commerçant.
